\begin{itemize}
    \item $A_0$ - The Sensible Object in Actuality (\gk{ἐνεργείᾳ} \gk{αἰσθητόν} / \textit{energeiai aisthēton})
    \begin{itemize}
        \item The sensible quality is actually exercising its power: the colour is visible, the sound audible. This is the \textit{first} actuality in the chain---the unmoved originator that initiates the entire perceptual event without itself being changed by the act of being perceived (\textit{DA} II.5, 417a-b).
        \item $M_0 \rightarrow M_1$: Perceptual Motion (\gk{αἴσθησις} / \textit{aisthēsis})
        \begin{itemize}
            \item The fulfillment of the sense organ's potential \textit{qua} capable of receiving the respective sensible form.
            \item Kinetic Roles:
            \begin{itemize}
                \item Unmoved Originator: $A_0$ - the sensible object in actuality
                \item Moved Mover: the sense faculty, which receives the form without the matter and is thereby altered (\textit{DA} II.12, 424a17-24)
                \item That Which is Moved: the sense organ and the perceiver as material substrate.
                \item Medium: air, water, or flesh---that through which the form is transmitted
                \item Time ($T_0 \rightarrow T_1$): the temporal interval over which the medium transmits and the organ receives. Nearly simultaneous for sight; measurably extended for sound.
            \end{itemize}
        \item This motion is named \textit{perception} because its terminus is $A_1$. The affective dimension is constitutive, not secondary: sensible form and affective valence arrive together as the world impresses itself on the soul.
        \item Destruction condition (426a27-b8): if the sensible exceeds the ratio of the organ, \hl{the motion acts on the organ \textit{as matter} rather than \textit{as sense}.} The ratio is destroyed; $A_1$ is never reached. The excess names no completed actuality and therefore constitutes no genuine perceptual motion---it is damage, not perception. \hl{Aristotle is describing a failed actualization of the sensory power, not the absence of motion tout court. There is real change (damage) in the organ, just not the right sort of change. This is no genuine perceptual motion in that there is no motion qua aisthēsis.}
        \end{itemize}
    \end{itemize}
\end{itemize}

\begin{itemize}
    \item $A_1$ - Completed Perception: Aisthēsis and Aisthēma (\gk{ἐνέργεια αἰσθήσεως} / \textit{energeia aistheseos})
    \begin{itemize}
        \item The sense has received the form without the matter. The ratio (\gk{λόγος} / \textit{lógos}) of the organ has been actualized within its proper range. This is a \textit{second actuality}: the organ possessed the capacity as first actuality (the sleeping geometer) and has now exercised it (the geometer proving). \hl{The actualization of the sensible object and that of the sensitive power are one and the same activity but different in being.}
        \item The \textit{aisthēma} is not merely a formal imprint; it is the perceptual act completed---both the sensible form and the soul's affective response (\gk{πάθος} / \textit{pathos}) are present as a unity. ``To perceive then is like bare asserting or thinking; but when the object is pleasant or painful, the soul makes a sort of affirmation or negation, and pursues or avoids the object'' (\textit{DA} III.7, 431a8--10). The perceptual faculty and the appetitive faculty are not two things joined together but one activity distinguished only in account: ``Both avoidance and appetite when actual are identical with this: the faculty of appetite and avoidance are not different, either from one another or from the faculty of sense-perception; but their being is different'' (431a13--14). The connection is analytic, not synthetic: wherever there is perception, there is pleasure and pain, and wherever these are, there is necessarily appetite (\textit{DA} II.3, 414b4--6). The world has disclosed itself to the perceiver in a determinate way---as both \textit{what it is} and \textit{how it matters}.
\inlinenote{In III.3, Aristotle uses aisthemata (in the plural) as what remain in the organs and can be re-activated, described as impression or picture left by the movement of perception. The dispositional element is textually backed by the analogy with knowledge vs. its exercise (III.2-3), where an acquired hexis can be re-activated.}
        \item $A_1$ names $M_0 \rightarrow M_1$. It also serves as the unmoved originator of $M_1 \rightarrow M_2$.
\inlinenote{II.5-6: the perceptual motion runs from object, through medium, to organ; the actuality of sense is the end-state. III.3: the aisthemata left behind become the basis for imagination; phantasia is defined as kinesis resulting from ``actual exercise of a power of sense,'' and that movement cannot exist apart from sensation.}
        \item $M_1 \rightarrow M_2$: \textit{Phantastic} Motion: From Residual Trace to \textit{Phantasma}
        \begin{itemize}
            \item After the sensible object withdraws, residual motion persists in the sensory apparatus (\textit{DA} III.3, 428b10--16; \textit{De Insomniis} 459a--461b). This is a single continuous process: the \textit{kinēsis} produced by perception persists, reverberates, and gradually settles into a determinate \textit{phantasma}. The residual motion already carries both dimensions of the original perceptual act---the formal imprint (\textit{aisthēma}) and the affective charge (\textit{pathos})---because it inherits the unity of the act from which it derives. There is no moment of fusion; there was never any separation.
            \item Kinetic Roles:
            \begin{itemize}
                \item Unmoved Originator: $A_1$---the completed perception (resonant \textit{aisthēma} and \textit{pathos}), which now functions as the source of a new motion.
                \item Moved Mover: \textit{phantasia} (\gk{φαντασία}) as the retentive faculty---sustains and transforms the impression, and in so doing is itself affected (carries the affection, \gk{πάθος} / \textit{pathos}). \hl{\textit{Phantasia} (the faculty/activity) and \textit{phantasma} (the image) are distinct. Imagination is not the image; it is the \textit{kinēsis} from sense ``in virtue of which'' images arise. Phantasia sustains, activates, and transforms the residual impression into a determinate image.}
                \item That Which is Moved: the psycho-somatic state of the animal, qualitatively changed by the retained impression. \hl{Psycho-somatic or psychophysical: Aristotle treats sensation, imagination, memory, and affections as ``enmattered accounts,'' as ``soul and body together.''}
                \item Time ($T_1 \rightarrow T_2$): the interval between the withdrawal of the sensible object and the settling of the residual trace into a determinate \textit{phantasma}. Brief for afterimages; extended for powerful impressions that settle slowly.
            \end{itemize}
        \item Named by its terminus: this is \textit{phantastic motion} because it terminates in $A_2$.
        \end{itemize}
    \end{itemize}
\end{itemize}

\begin{itemize}
      \item $A_2$ --- The \textit{Phantasma} Proper (\gk{φάντασμα})
      \begin{itemize}
          \item The determinate image: the residual motion has settled into a
          dual-aspect unity of formal imprint (\textit{aisthēma}) and affective
          charge (\textit{pathos}) in which the sensible form of the object and the
          soul's dispositional orientation toward it achieve determinate unity.
          The \textit{phantasma} is:
          \begin{itemize}
              \item Stable and recallable---available as an object for further
              cognitive operations.
              \item Formally similar to the original perception but operating in its
              absence (\textit{DA} III.3, 428b10--16: ``imagination is held to be a
              movement \ldots\ necessarily similar in character to the sensation
              itself'').
              \item The hinge between perception and thought---``the soul never
              thinks without an image'' (\textit{DA} III.7, 431b2).
          \end{itemize}
          \item $A_2$ names $M_1 \rightarrow M_2$. It serves as the unmoved originator of
          $M_2 \rightarrow M_3$. This is the pivotal actuality: everything prior
          leads to it; everything subsequent operates upon it.

          \item $M_2 \rightarrow M_3$: \textbf{Noetic / Doxastic / Mnemonic / Anticipatory
          Motion: Cognitive Engagement with the \textit{Phantasma}}
          \begin{itemize}
              \item The \textit{phantasma} is taken up by \textit{nous}, the doxastic
              faculty, memory, or deliberative \textit{phantasia}. These are
              \textit{parallel modes} distinguished by intentional direction,
              temporal orientation, and relationship to truth, not by position in the
              causal chain:
\inlinenote{Aristotle distinguishes sensitive and deliberative phantasia (433b29--434a10) but never explicitly says they are parallel orientations of the same apparatus. The unification is implicit in the texts. Likewise, he never says that temporal awareness is what makes possible deliberative phantasia; the capacity to calculate and the capacity to perceive time seem to be co-implicated. Bowin argues that deliberative phantasia actively manipulates images to produce determinate temporal representations---treating it as a function of deliberative phantasia, not its condition.}
              \begin{enumerate}
                  \item \textbf{Intellection}: the \textit{phantasma} addressed
                  \textit{as bearing a universal}---\textit{nous} abstracts the form
                  from the image. Temporal orientation: \textit{atemporal}. The
                  universal, once grasped, is not indexed to past, present, or
                  future but holds across all instances.

                  \item \textbf{Doxa} (\gk{δόξα}): the \textit{phantasma} addressed
                  \textit{as warranting propositional assertion}---the soul commits
                  itself to a truth-claim about what is the case. Doxa is the
                  \textit{logos}-dependent, assertoric operation on the
                  \textit{phantasma} that phantasia itself is not: ``imagination does
                  not involve opinion based on inference, though opinion involves
                  imagination'' (\textit{DA} III.11, 434a10--11). It requires three
                  conditions that phantasia lacks: belief (\gk{πίστις} /
                  \textit{pistis}), conviction (\gk{πεπεῖσθαι} /
                  \textit{pepeisthai}), and discourse of reason (\gk{λόγος}): ``opinion
                  involves belief, belief involves being convinced, and being
                  convinced involves reason; but some of the brutes have
                  imagination, without discourse of reason'' (\textit{DA} III.3,
                  428a19--24). Accordingly, doxa is found only in rational animals.

                  Doxa differs from the other modes in two decisive respects:

                  \begin{itemize}
                      \item \textit{Involuntariness}: ``imagining lies within our
                      power whenever we wish \ldots\ but in forming opinions we are
                      not free: we must either be holding what is true or what is
                      false'' (\textit{DA} III.3, 427b17--21). Phantasia is
                      \gk{ἐφ' ἡμῖν}; doxa is \gk{οὐκ ἐφ' ἡμῖν}. One can summon
                      any image at will, but one cannot choose what one takes to be
                      the case. Doxa arises from the way things appear under the
                      sway of phantasia---the taking-as function---but once it
                      arises, the soul is bound by its content. The involuntariness
                      is grounded externally: doxa is answerable to how things are,
                      not to how we might wish them to appear.
                      \item \textit{Independence from the phantasma's own testimony}:
                      doxa and phantasia can deliver contradictory content
                      simultaneously. ``We imagine the sun to be a foot in diameter
                      though we are convinced (\gk{πεπεῖσθαι}) that it is larger
                      than the inhabited earth'' (\textit{DA} III.3, 428b2--4). The
                      \textit{phantasma} is not corrected by doxa---it persists
                      unchanged---but doxa is not compelled by the
                      \textit{phantasma}'s apparent testimony either. Doxa is thus a
                      distinct assertoric act that takes the \textit{phantasma} as
                      material but can override or diverge from its content. This is
                      why ``a true opinion becomes false only when the fact alters
                      without being noticed; imagination is therefore neither any one
                      of the states enumerated, nor compounded out of them''
                      (428b4--10).
                  \end{itemize}

                  The formation of doxa is subject to a \textit{gatekeeping
                  condition}: whether the corrective rational faculty intercepts the
                  phantasma before doxastic commitment occurs. When the controlling
                  faculty operates normally, the phantasma may be corrected (as in
                  the sun case: we see ``one foot'' but do not believe it). When the
                  corrective faculty is disabled---by sleep, illness, passion, or
                  distance---the phantasma passes unchallenged into doxa: ``when the
                  controlling sense is not disabled \ldots\ we are not deceived; but
                  when it is disabled, as happens in sleep and illness, then what
                  appears is taken to be what is the case'' (\textit{De Insomniis}
                  460b3--16; cf.\ \textit{DA} III.3, 428b25--429a2: phantasia is
                  ``especially'' erroneous ``when the object of perception is far
                  off''). The speed of the transition can be near-instantaneous:
                  ``when we take something to be fearful (\gk{δοξάσωμεν}),
                  emotion is immediately produced (\gk{εὐθὺς πάσχομεν})''
                  (\textit{DA} III.3, 427b21--24).

                  Temporal orientation: primarily \textit{present-situational}---doxa
                  asserts ``this is so'' or ``this is not so'' about what the
                  phantasma discloses---though doxa may concern past or future states
                  of affairs. Its temporal character is assertoric rather than
                  oriented: it \textit{commits} without privileging a temporal
                  direction.

\inlinenote{The relationship between doxa and the practical syllogism is structurally significant. In \textit{MA} 701a7--25, the practical syllogism proceeds from a universal premise (which may be a settled doxa or epistēmē: ``all sweet things are to be tasted'') conjoined with a particular perception or phantasma (``this is sweet'') to issue in action. Doxa can thus function both as an \textit{output} of the current perceptual event (the particular judgment ``this is dangerous'') and as a pre-existing \textit{hexis} (a settled universal belief) that informs how subsequent phantasmata are processed. In the latter case, the universal doxa operates as an additional unmoved originator alongside $A_2$: the repertoire of settled beliefs shapes how the phantasma is taken up in $M_2 \rightarrow M_3$.}

                  \item \textbf{Memory}: the \textit{phantasma} addressed \textit{as
                  a likeness of something past}---the past is marked as past and
                  felt as past (\textit{On Memory} 449b24--25). Temporal orientation:
                  \textit{retrospective}. Memory is constitutively temporal; its
                  object \textit{is} the past \textit{qua} past.

                  \item \textbf{Deliberative / Anticipatory \textit{Phantasia}}
                  (\gk{βουλευτικὴ φαντασία}): \textit{several} \textit{phantasmata}
                  synthesized into a unity and addressed \textit{as composable into
                  unrealized possibilities}, measured by a single standard
                  (\textit{DA} III.11, 434a5--10: ``whether this or that shall be
                  enacted is already a task requiring calculation; and there must be
                  a single standard to measure by, for that is pursued which is
                  greater. It follows that what acts in this way must be able to make
                  a unity out of several images''). Temporal orientation:
                  \textit{prospective}. This mode operates in two interlocking
                  sub-functions:
                  \begin{itemize}
                      \item \textit{Projective-synthetic}: the \textit{phantastic}
                      faculty recombines elements drawn from the available repertoire
                      of \textit{phantasmata} into images of states not yet
                      actual---possible futures, anticipated outcomes, unrealized
                      configurations. This is the productive capacity that goes beyond
                      retention: \textit{phantasia} here constructs rather than
                      merely preserves.
                      \item \textit{Calculative-practical}: the projected
                      possibilities are measured against one another by a single
                      standard (\gk{ἑνί τινι δεῖ μετρεῖν}) to determine which
                      course of action shall be enacted. This is where
                      \textit{phantasia} converges with practical thought
                      (\gk{διάνοια πρακτική}): the projected image is evaluated not
                      merely as possible but as \textit{to-be-pursued} or
                      \textit{to-be-avoided}.
                  \end{itemize}
                  This dual structure is what distinguishes rational action from
                  appetitive reaction. Non-rational animals act from sensitive
                  \textit{phantasia} operating on a single present
                  \textit{phantasma}---``I want to drink, says appetite; this is
                  drink, says sense or imagination: straightaway I drink''
                  (\textit{MA} 701a32--33). Rational animals, because they possess
                  \gk{βουλευτικὴ φαντασία}, can hold back from what is immediately
                  present by projecting what is future---and this capacity requires
                  a \textit{sense of time}: ``while thought bids us hold back because
                  of what is future, desire is influenced by what is just at hand: a
                  pleasant object which is just at hand presents itself as both
                  pleasant and good, without condition in either case, because of
                  want of foresight into what is farther away in time'' (\textit{DA}
                  III.10, 433b5--10). The sense of time is co-implicated with
                  deliberative \textit{phantasia}: the capacity to calculate and the
                  capacity to perceive time cannot be separated.

\inlinenote{This is why thought can override desire when considering future implications. Pursuit of the good still drives all movement, but in a calculative, prospective fashion.}

              \end{enumerate}

              \item These modes may operate simultaneously or in succession. Memory
              is not a separate motion-stage but a specific orientation of the same
              apparatus. Likewise, anticipatory \textit{phantasia} is not a separate
              faculty but the \textit{prospective} orientation of the same
              \textit{phantastic} capacity that, turned backward, yields memory and,
              turned inward, furnishes \textit{nous} with its images. Doxa, for its
              part, is the \textit{assertoric} engagement with the
              \textit{phantasma}---the soul's involuntary commitment to the
              appearance as true or false---which may accompany or inflect any of
              the other modes. One may remember \textit{and} hold a belief about what
              is remembered; one may deliberate \textit{and} form opinions about the
              projected outcomes. Doxa is not a separate temporal orientation but a
              truth-committing layer that can supervene on any of the three
              orientations (atemporal, retrospective, prospective).

              \item Kinetic Roles:
              \begin{itemize}
                  \item Unmoved Originator: $A_2$---the \textit{phantasma} (or, in
                  the case of deliberative \textit{phantasia}, the available
                  \textit{repertoire} of \textit{phantasmata}) as object of thought
                  (\textit{DA} III.7, 431b2). For doxa specifically, settled
                  universal beliefs (\textit{doxai} as \textit{hexeis}) may also
                  function as additional unmoved originators, shaping how the
                  \textit{phantasma} is taken up (cf.\ the universal premise in the
                  practical syllogism, \textit{MA} 701a7--25).
                  \item Moved Mover: \textit{nous} (\gk{νοῦς}) (for intellection),
                  the doxastic faculty (for opinion-formation: \gk{λόγος} operating
                  through \gk{πίστις} and \gk{πεπεῖσθαι}; \textit{DA} III.3,
                  428a19--24), the mnemonic faculty (for memory),
                  \gk{βουλευτικὴ φαντασία} (for deliberation and
                  anticipation)---each activated by and operating upon the
                  \textit{phantasma}.
                  \item That Which is Moved: the knower / believer / rememberer /
                  deliberator / anticipator---the animal is cognitively transformed.
                  \item Time ($T_2 \rightarrow T_3$): variable by mode. Intellection
                  of universals can be instantaneous---the universal is atemporal in
                  character, grasped in a single act rather than extended over
                  duration. Doxa-formation varies: simple perceptual doxa (``this is
                  white,'' ``this is dangerous'') may be near-instantaneous, its speed
                  indicated by Aristotle's \gk{εὐθύς} (``immediately,'' \textit{DA}
                  III.3, 427b23); inferential doxa, which involves combining premises
                  through \textit{logos}, extends over the duration of the reasoning
                  process. Memory is constitutively temporal: its object \textit{is}
                  the past \textit{qua} past, and the temporal interval between the
                  original perception and the present act of recall is internal to
                  what memory discloses. Deliberative / anticipatory
                  \textit{phantasia} is constitutively \textit{prospective}: it
                  extends over the weighing of alternatives projected into the
                  future, and its temporal horizon is bounded by the ``sense of
                  time'' (\textit{DA} III.10, 433b7) that enables foresight into
                  what is farther away. Time here is not merely the measure of the
                  motion but the medium within which the four modes differentiate
                  themselves by intentional direction: beyond duration (intellection),
                  assertoric commitment (doxa), backward (memory), and forward
                  (anticipation).
              \end{itemize}

              \item Named by its terminus: this is
              \textit{thinking / opining / remembering / deliberating / anticipating}
              because it terminates in $A_3$.
          \end{itemize}
      \end{itemize}
  \end{itemize}

  \begin{itemize}
      \item $A_3$ --- Completed Noetic / Doxastic / Mnemonic / Anticipatory Actuality
      \begin{itemize}
          \item The judgment is formed. The universal is grasped, the opinion is
          held, the past is recalled and recognized as past, the deliberation has
          reached its conclusion, or the anticipated possibility has been
          constructed, evaluated, and either adopted or rejected. The
          \textit{phantasma}---or, in the case of deliberative \textit{phantasia},
          the synthesized unity of several \textit{phantasmata}---has been fully
          processed, refined in both its formal content and its felt significance.

          \item In the case of doxa, $A_3$ is a determinate belief: the soul now
          holds that something \textit{is} or \textit{is not} the case. This
          completed opinion differs from the completed acts of the other modes in
          its involuntary bindingness---once formed, the subject is not free to
          discard it at will (\textit{DA} III.3, 427b17--21)---and in its
          distinctive affective consequence. Where the \textit{phantasma} at $A_2$
          carries latent affective valence (\textit{pathos} as dispositional
          orientation), doxa \textit{activates} that valence into a condition
          capable of producing higher-order emotion: ``when we merely imagine
          something fearful or encouraging we remain unaffected
          (\gk{ἀπαθῶς ἔχομεν}), but when we take something to be fearful
          (\gk{δοξάσωμεν}), emotion is immediately produced
          (\gk{εὐθὺς πάσχομεν})'' (\textit{DA} III.3, 427b21--24). The
          \textit{phantasma}'s latent \textit{pathos} is necessary but not
          sufficient; it is doxa's truth-commitment that converts dispositional
          orientation into the mobilized affective state from which higher-order
          emotions (fear, anger, desire, etc.) can arise. The analysis of emotion
          as a distinct moment in the chain is reserved for subsequent treatment.

\inlinenote{The distinction between \textit{pathos} (affective valence present from $A_1$ onward) and \textit{emotion} (the higher-order affective states---fear, anger, pity, etc.---that arise only after doxastic commitment) is structurally load-bearing. Pathos is the raw good/bad, pleasant/painful orientation impressed with the formal trace at the moment of perception and inherited by the phantasma. Emotion is the determinate mobilization of the soul that follows from \textit{taking-as-true}: the soul is not merely disposed toward avoidance (pathos) but is \textit{afraid} (emotion). The passage at 427b21--24 names precisely this transition: \gk{ἀπαθῶς ἔχομεν} under bare phantasia; \gk{εὐθὺς πάσχομεν} under doxa. The emotion moment will be developed in a subsequent revision of the chain.}

          \item What emerges is the apprehension of a realizable good (or evil to
          be avoided): the object of desire that initiates the final motion.
          Judgment and affective evaluation have converged---either explicitly
          through deliberation or compressed into kinesthetic schemata of
          habituated practice. In the case of deliberative \textit{phantasia}, the
          ``realizable good'' is not merely recognized in what is present but
          \textit{projected} into an anticipated future and selected from among
          competing possibilities by a single standard of measure. In the case of
          doxa, the realizable good (or evil) is disclosed through the soul's
          involuntary commitment to the appearance as true: the object is not
          merely imagined as threatening but \textit{taken to be} threatening, and
          this taking-as-true is what gives it motivational force. Speculative doxa,
          however, need not issue in action at all: ``practical thought [which
          calculates means to an end] differs from speculative thought'' in that the
          latter ``moves nothing'' (\textit{DA} III.10, 433a13--15). Where doxa is
          speculative---where the soul holds a belief about what is the case without
          that belief disclosing a realizable good or evil---$A_3$ does not serve as
          unmoved originator of $M_3 \rightarrow M_4$. The chain terminates at $A_3$
          for that mode.

          \item $A_3$ names $M_2 \rightarrow M_3$. It serves as the unmoved
          originator of $M_3 \rightarrow M_4$---but only mediately, through the
          realizable good it discloses. The good that moves appetite may be
          immediately present (sensitive \textit{phantasia}: the drink before me)
          or anticipatorily projected (deliberative \textit{phantasia}: the
          course of action whose outcome I foresee) or involuntarily taken to be
          the case (doxa: the thing believed dangerous). In either case, it is $A_3$
          as completed actuality---not the raw \textit{phantasma}---that functions
          as the unmoved originator of desire.
\inlinenote{433b10--18 describes the object of desire as that which moves without being moved, and desire as that which moves and is moved.}
          \item $M_3 \rightarrow M_4$: \textbf{Orektikon Motion: Desire Converging
          in Action}
          \begin{itemize}
              \item Judgment and desire converge in overt action, drawing on the
              full kinetic-affective history of the \textit{phantasma}. Where
              deliberative \textit{phantasia} has operated, the desire is informed
              by foresight: the animal acts not merely from what is at hand but
              toward what has been projected and evaluated as the greater good.
              Where only sensitive \textit{phantasia} is at work, appetite responds
              to the immediately present \textit{phantasma} without foresight---
              ``because of want of foresight into what is farther away in time''
              (\textit{DA} III.10, 433b9--10). Where doxa has operated, the
              desire is informed by the soul's truth-commitment: the animal acts
              not merely from how things appear but from how they are
              \textit{taken to be}. The affective force of doxa-mediated desire
              is immediate (\gk{εὐθύς}, \textit{DA} III.3, 427b23) and
              involuntary (\gk{οὐκ ἐφ' ἡμῖν}, 427b20)---the soul that believes
              something dangerous is already mobilized toward avoidance before
              deliberation can intervene.
              \item Kinetic Roles:
              \begin{itemize}
                  \item Unmoved Originator: the realizable good as apprehended in
                  $A_3$ (\textit{DA} III.10, 433b10--18).
                  \item Moved Mover: appetite / desire (\gk{ὄρεξις} /
                  \textit{orexis})---``that which at once moves and is moved''
                  (433b17--18); ``appetite in the sense of actual appetite
                  \textit{is} a kind of movement'' (433b18).
                  \item That Which is Moved: the animal via the bodily
                  instrument---the ball-and-socket joint where ``pushing and
                  pulling'' execute the action (433b21--25).
                  \item Time ($T_3 \rightarrow T_4$): the interval between the
                  formation of desire and the completion of the motor act.
              \end{itemize}
              \item Named by its terminus: this is \textit{appetitive action}
              because it terminates in $A_4$.

\inlinenote{the source of all animal movement. MA CITE}

          \end{itemize}
      \end{itemize}
  \end{itemize}

  \begin{itemize}
      \item $A_4$ --- Completed Action (\gk{πρᾶξις})
      \begin{itemize}
          \item The perceptual-\textit{phantasmatic}-motor loop is closed. The
          ensouled body has engaged the moving world. The action draws on the full
          kinetic-affective history of the \textit{phantasma}---from the original
          perceptual encounter through its retention, cognitive
          reactivation, and appetitive appropriation.
          \item $A_4$ names $M_3 \rightarrow M_4$. The completed action changes the
          world---it produces new sensible conditions that can serve as $A_0$ for a
          new perceptual cycle. The chain is not linear but recursive.
      \end{itemize}
  \end{itemize}
